% Nouveau EGA V (cote 25): archive cote 25, site page 86 (PDF page 87).
% typescript. Draft compared: none.
% AI-transcribed by Claude Opus 5.5 (claude-opus-5-5, adaptive thinking, effort high), Message Batches API, batch ega25b (msgbatch_01A7n5YuXubUPNK1zsepi8RG), 2026-11-13; not yet checked by eye.
\scanpage{25/86}
\mspage{\hand{\unsure{80}}}

p$_{\unsure{\xout{C}}}$ précédent, \deleted{qu'il sera commode de noter également de la même façon}
\deleted{p$_C$ , pour indiquer sa dépendance de C.} De cette façon, les propriétés de
la projection conique  restreinte  sont ramenées \deleted{dans une large mesure} à
celles de la projection conique étendue, qui a été  étudiée systématique-%
ment par ailleurs, ou est sensée l'avoir été (cf n\textsuperscript{o} 10 et n\textsuperscript{o} 12). La
question\xout{s} principale\xout{s} qui se pose est, lorsque  S=Spec(k), quelles sont
les propriétés de la p\unsure{ro}jection conique  de X  quand on prend C générique
dans Grass$^{m+1}$(\unsure{E|P}) (ce qui exige  qu'on fasse un changement de base
k$\rightarrow$k(\hand{$\unsure{\eta|y}$}) , i.e.  C est \inserted{alors} en fait une sous-variété linéaire de X$_{k(\hand{\eta})}$ ); des
\unsure{ar}guments  standart, déjà répétés à satiété, permettront d'en conclure
les propriétés analogues pour les projections coniques correspondant à \xout{de}
des points de Grass$^{m+1}$(P)  appartenant à un ouvert non vide de ladite
grassmanienne, et enfih, lorsque k est infini, on conclut à l'existence
d'une(en fait une infinité\hand{)} \inserted{de} \xout{de} C  \uline{définie\xout{s} sur} k , I.e.  une sous-variété
linéaire de P lui-même (sans changement de corps de base), donnant lieu
à une projection conique ayant les propriétés en question. Il y aurait
\margin{|} lieu de bloquer ce genre d'explications générales avec celles de la même
\margin{|} eau données dans les n\textsuperscript{os} 4 et 7, et que nous avons d'ailleurs déjà
\margin{|} utilisées plus ou moins implicitement,  par exemple au n\textsuperscript{o}13 \xout{\unsure{(dém. de}}
\xout{\unsure{13.3}} (cf 13.4. c). Il y a lieu également, à ce propos, d'examiner les
propriétés relatives à un faisceau F sur X,  en prenant son image inverse
$F^{(m)}_{Q(C)}$  sur  $\underline{\underline{X}}^{(m)}_{Q(C)}$ . Il faut  d'ailleurs, dans la situation précise
envisagée ici, d es notations plus simples, je propose  $\widetilde{X(C)}$ et $\widetilde{F(C)}$ ,
et simplement $\widetilde{X}$ et $\widetilde{F}$  quand aucune confusion n\unsure{'}en résulte  (attention, le
F \deleted{précédent} \inserted{\unsure{ici}} n'est pas  celui du début du n\textsuperscript{o}).

Grosso modo, et lorsque  disons on suppose  f une immersion, les
propriétés de la projection conique générique sont assez différentes,

\mspage{\hand{86}}
